\documentclass[12pt,a4paper]{article}
\usepackage[utf8]{inputenc}
\usepackage[T1]{fontenc}
\usepackage[brazilian]{babel}
\usepackage{geometry}\geometry{margin=2cm}
\usepackage{booktabs}\usepackage{tabularx}\usepackage{enumitem}\usepackage{hyperref}\usepackage{xcolor}\usepackage{longtable}\usepackage{array}
\setlength{\parskip}{0.5em}\setlength{\parindent}{0pt}
\definecolor{decidido}{RGB}{20,110,80}\definecolor{pendente}{RGB}{176,98,10}\definecolor{futuro}{RGB}{80,80,80}
\newcommand{\statusDecidido}{\textcolor{decidido}{\textbf{DECIDIDO}}}\newcommand{\statusPendente}{\textcolor{pendente}{\textbf{PENDENTE}}}\newcommand{\statusFuturo}{\textcolor{futuro}{\textbf{FUTURO}}}
\title{\textbf{Núcleo de Situação de Saúde}\\Perguntas, decisões e escopo epidemiológico}\author{Documento para decisão conjunta da equipe}\date{Outubro de 2026}
\begin{document}\maketitle\tableofcontents\newpage
\section{Objetivo deste documento}
Este documento registra duas coisas diferentes e igualmente importantes:
\begin{enumerate}[leftmargin=*]\item as perguntas de vigilância, gestão, planejamento e comunicação que podem orientar a evolução futura do NSS;\item as decisões já tomadas para a \textbf{NSS 1.0 / Frontend V2}, evitando que pipeline, Java e frontend interpretem o escopo de maneiras diferentes.\end{enumerate}
A disponibilidade de um campo ou dataset não é motivo suficiente para incorporá-lo. Primeiro define-se a pergunta; depois são escolhidos fonte, campos, regra de qualidade, agregação e visualização.
\section{Estado das decisões - NSS 1.0}
\begin{longtable}{p{4.4cm}p{2.4cm}p{8.2cm}}\toprule\textbf{Tema} & \textbf{Estado} & \textbf{Decisão} \\\midrule\endfirsthead\toprule\textbf{Tema} & \textbf{Estado} & \textbf{Decisão} \\\midrule\endhead
Escopo de doenças & \statusDecidido & Dengue, Febre Maculosa e Toxoplasmose. A base/código exato de Toxoplasmose ainda precisa de validação epidemiológica antes de ser anunciada como disponível. \\
Métrica exibida & \statusDecidido & Usar o termo \textbf{notificações} enquanto a regra de casos confirmados por agravo não estiver formalmente aprovada. \\
Municípios & \statusDecidido & Campos dos Goytacazes, São João da Barra, Macaé e Itaperuna. \\
Perspectiva geográfica & \statusDecidido & Notificação, não residência. \\
Detalhamento em Campos & \statusDecidido & Município -> distrito -> bairro/localidade da unidade notificadora. \\
Demais municípios & \statusDecidido & Somente nível municipal na NSS 1.0. \\
Tempo & \statusDecidido & Ano e mês derivados de \texttt{DT\_NOTIFIC}. \\
Semana epidemiológica & \statusFuturo & Mantida em Silver/Gold Ad Hoc para auditoria, fora da Serving/UI V1.0. \\
Sexo & \statusDecidido & Todos, Masculino e Feminino. ``Todos'' significa ausência de filtro e preserva ignorados. \\
Faixa etária & \statusDecidido & <1; 1--4; 5--9; 10--14; 15--19; 20--39; 40--59; 60--64; 65--69; 70--74; 75--79; 80+. \\
Origem da idade & \statusDecidido & Derivar a faixa de \texttt{NU\_IDADE\_N}; \texttt{ANO\_NASC} não é suficiente para o contrato Serving. \\
Multi-faixa & \statusFuturo & Não entra na V1.0. O banco/Java deve permitir evolução futura sem pré-calcular todas as combinações. \\
Bairro de residência & \statusFuturo & Depende de parceria institucional com o CIEVS e fonte autorizada. \\
CEP de residência & \statusDecidido & Não é necessário para a finalidade atual e não deve ser solicitado para esse objetivo. \\
Identificadores pessoais & \statusDecidido & Nome, CPF, CNS, telefone, endereço completo e equivalentes não são necessários à finalidade atual. \\
CIDAC & \statusDecidido & Referência conceitual para organização distrito -> bairro/localidade, não fonte canônica de dados/geometria. \\
GeoJSON & \statusDecidido & Gerado por Python para uso do frontend. \\
UBS no mapa & \statusDecidido & Localização pública client-side; não precisa participar do grão analítico da Gold Serving. \\
Casos confirmados & \statusPendente & Definir por agravo, a partir de dicionário oficial e validação epidemiológica, quais classificações entram no indicador de caso. \\
Toxoplasmose & \statusPendente & Confirmar qual base/código do SINAN representa o indicador desejado; não somar categorias distintas automaticamente. \\
Pequenas contagens & \statusFuturo & Política institucional necessária antes de abertura pública irrestrita e antes de uso de dados restritos. Não bloqueia a primeira versão autenticada com fonte pública. \\
Estado RJ completo & \statusFuturo & Versão posterior. \\
Região Sudeste & \statusFuturo & Somente após cobertura dos estados desejados. \\
\bottomrule\end{longtable}
\section{Contexto atual do Núcleo}
O NSS está em fase inicial de integração. A iniciativa pretende apoiar saúde humana e animal por meio de monitoramento epidemiológico, planejamento de contingência, apoio a órgãos públicos e comunicação de informação de saúde.
A fonte atualmente incorporada à pipeline é o \textbf{SINAN via PySUS}. O front-end V1 já demonstrou navegação Região -> Estado -> Município com dados de demonstração. A NSS 1.0 passa a integrar dados reais, autenticação, PostgreSQL, Java e Frontend V2.
A futura parceria com o CIEVS é importante para dados territoriais de residência que não estão disponíveis no dataset público atualmente utilizado. A arquitetura deve deixar rota futura para essa integração sem tornar a NSS 1.0 dependente dela.
\section{O que a pipeline atual já faz}
A pipeline do repositório \texttt{Zadoque/NSS-pipeline} segue arquitetura Medallion e já evoluiu para carregar PostgreSQL:
\begin{center}\texttt{PySUS/SINAN $\rightarrow$ Bronze $\rightarrow$ Silver $\rightarrow$ Gold $\rightarrow$ PostgreSQL}\end{center}
\begin{itemize}[leftmargin=*]\item \textbf{Bronze}: entrada e metadados;\item \textbf{Silver}: valida, normaliza, filtra e deduplica;\item \textbf{Gold Ad Hoc}: flexível para inspeção/exploração;\item \textbf{Gold Serving}: contrato fixo da NSS 1.0;\item \textbf{PostgreSQL}: camada analítica consumida pelo Java.\end{itemize}
A auditoria de outubro de 2026 identificou correções necessárias antes da territorialização por unidade notificadora, especialmente na ingestão/paginação do CNES, semana epidemiológica e derivação da faixa etária.
\subsection{Campos centrais para a NSS 1.0}
\begin{longtable}{p{4.2cm}p{3.3cm}p{7.1cm}}\toprule\textbf{Conceito} & \textbf{Campo/fonte} & \textbf{Uso na V1.0} \\\midrule\endfirsthead\toprule\textbf{Conceito} & \textbf{Campo/fonte} & \textbf{Uso na V1.0} \\\midrule\endhead
Data de notificação & \texttt{DT\_NOTIFIC} & Define ano e mês. \\
Município da notificação & \texttt{ID\_MUNICIP} & Nível municipal. \\
UF de notificação & \texttt{SG\_UF\_NOT} & Validação/recorte. \\
Unidade notificadora & \texttt{ID\_UNIDADE} & Join CNES para território intramunicipal. \\
Sexo & \texttt{CS\_SEXO} & Filtro M/F; ignorados preservados. \\
Idade & \texttt{NU\_IDADE\_N} & Derivação das 12 faixas + NI. \\
Classificação final & \texttt{CLASSI\_FIN} & Futura regra formal de casos confirmados. \\
Identificador da notificação & \texttt{NU\_NOTIFIC} e/ou chave validada & Deduplicação; unicidade auditada por dataset. \\
CNES & estabelecimento & Território da unidade notificadora. \\
\bottomrule\end{longtable}
\section{Dados públicos versus futura fonte CIEVS}
O modelo SINAN possui campos de residência, mas a disponibilidade efetiva depende do dataset disseminado. O Parquet público atualmente utilizado não deve ser tratado como se contivesse bairro de residência quando esse campo não está presente.
A NSS 1.0 usa a localização da unidade notificadora, disponível a partir de \texttt{ID\_UNIDADE} + CNES. A perspectiva por residência será funcionalidade separada quando houver fonte institucional autorizada.
\subsection{Princípio de minimização para a futura parceria}
Para residência por território, a necessidade provável inclui município/distrito/bairro de residência, sexo, idade, período, classificação e chave técnica pseudonimizada. Não há justificativa atual para pedir nome, CPF, CNS, telefone, logradouro, número ou CEP.
\section{Perguntas candidatas para evolução futura}
As perguntas abaixo permanecem como backlog; classificação futura não significa descarte.
\subsection{Situação epidemiológica atual}
\begin{enumerate}[label=\textbf{E\arabic*.},leftmargin=*]\item Quantas notificações e, quando validado, quantos casos confirmados existem por período e território?\item Como os números evoluem semana a semana e mês a mês?\item Quais municípios apresentam aumento recente?\item Quais doenças crescem simultaneamente?\item Qual é a distribuição por sexo e faixa etária?\item Qual é a distribuição por classificação final/evolução?\item Quantos permanecem sem classificação conclusiva?\item Há concentração espacial em municípios vizinhos?\end{enumerate}
Na NSS 1.0, E1 e E5 são priorizadas no recorte aprovado.
\subsection{Indicadores e comparação justa entre territórios}
\begin{enumerate}[label=\textbf{I\arabic*.},leftmargin=*]\item Qual é a incidência por 10 mil/100 mil habitantes?\item Como o município se compara ao próprio histórico?\item Como se compara a municípios semelhantes?\item Qual a tendência móvel?\item Quais indicadores caracterizam possível surto?\end{enumerate}
\subsection{Planejamento e contingência}
\begin{enumerate}[label=\textbf{C\arabic*.},leftmargin=*]\item Em quais territórios crescimento sugere plano de contingência?\item Onde estão serviços capazes de atender a demanda?\item Quais municípios precisam de apoio logístico?\item Há capacidade assistencial próxima?\item Como priorizar insumos/equipes?\end{enumerate}
\subsection{Vacinação e prevenção}
\begin{enumerate}[label=\textbf{V\arabic*.},leftmargin=*]\item Cobertura vacinal adequada?\item Áreas com baixa cobertura e aumento de doença?\item Quais faixas/territórios priorizar?\end{enumerate}
\subsection{Mortalidade e gravidade}
\begin{enumerate}[label=\textbf{M\arabic*.},leftmargin=*]\item Houve aumento de óbitos?\item Relação notificações/internações/óbitos?\item Quais doenças têm evolução desfavorável?\end{enumerate}
\subsection{Saúde ambiental e determinantes}
\begin{enumerate}[label=\textbf{A\arabic*.},leftmargin=*]\item Associação com qualidade da água?\item Territórios com risco sanitário ambiental?\item Quais indicadores ambientais exibir?\end{enumerate}
\subsection{Saúde animal e abordagem integrada}
\begin{enumerate}[label=\textbf{Z\arabic*.},leftmargin=*]\item Quais zoonoses priorizar?\item Quais perguntas exigem dados veterinários?\item Quais órgãos/sistemas integrar?\item Correlacionar eventos animais/humanos?\item Quais dados podem ser compartilhados entre vigilâncias?\end{enumerate}
\section{Perguntas de qualidade e auditoria}
\begin{enumerate}[label=\textbf{Q\arabic*.},leftmargin=*]\item Qual batch/data da última atualização?\item Quantos registros entraram em Bronze e permaneceram em Silver?\item Qual soma foi publicada na Gold Serving?\item A mesma soma está no PostgreSQL?\item Quantas notificações não mapearam CNES/território?\item O indicador é notificação ou caso confirmado?\item Qual versão de contrato/processamento o originou?\end{enumerate}
Para NSS 1.0, Grafana prioriza Bronze, Silver, Gold e PostgreSQL com reconciliação mínima.
\section{Decisões estruturais consolidadas}
\subsection{Definição de ``caso''}\statusPendente. A regra deve ser específica por agravo e validada por dicionário oficial/equipe epidemiológica. Enquanto isso, usar \textbf{notificações}.
\subsection{Ano e período}\statusDecidido. Ano/mês de \texttt{DT\_NOTIFIC}; semana fica fora do filtro oficial V1.0.
\subsection{Granularidade geográfica}\statusDecidido. Perspectiva de notificação; quatro municípios em nível municipal e somente Campos com drill-down por distrito/bairro da unidade notificadora.
\subsection{Faixa etária}\statusDecidido. Derivar de \texttt{NU\_IDADE\_N}; produzir 12 faixas + NI. Não pré-calcular combinações de multi-faixa.
\subsection{Privacidade e pequenas contagens}\statusFuturo para publicação pública. A NSS 1.0 é autenticada e usa fonte pública agregada; regra institucional será necessária antes de abertura pública e dados restritos.
\section{Matriz de prioridade da NSS 1.0}
\begin{longtable}{p{6.1cm}p{2.4cm}p{2.2cm}p{4.1cm}}\toprule\textbf{Pergunta/capacidade} & \textbf{Prioridade} & \textbf{Release} & \textbf{Fonte} \\\midrule
Notificações por quatro municípios/ano/mês & Alta & NSS 1.0 & SINAN/PySUS \\
Filtro por sexo & Alta & NSS 1.0 & \texttt{CS\_SEXO} \\
Filtro por 12 faixas etárias & Alta & NSS 1.0 & \texttt{NU\_IDADE\_N} \\
Distrito/bairro da unidade em Campos & Alta & NSS 1.0 & \texttt{ID\_UNIDADE}+CNES+referência territorial \\
Login/acesso autenticado & Alta & NSS 1.0 & Java/Spring Security \\
Auditoria Bronze/Silver/Gold/PostgreSQL & Alta & NSS 1.0 mínima & Pipeline/PostgreSQL/Grafana \\
Bairro de residência & Posterior & V1.1+ & CIEVS/restrita \\
Multi-faixa & Posterior & V1.1+ & SQL sobre Serving \\
Todos os municípios RJ & Posterior & V2+ & SINAN/PySUS \\
Incidência populacional & Posterior & V2+ & SINAN+IBGE \\
Saúde animal & Posterior & V2+ & fontes a definir \\
\bottomrule\end{longtable}
\section{Critérios para futuras decisões}
Para cada pergunta aprovada, registrar definição do indicador, evento contado, período, perspectiva geográfica, fonte/campos mínimos, regra de qualidade, política de privacidade, periodicidade, visualização e responsável pela validação epidemiológica.
\section{Referências técnicas}
\begin{itemize}[leftmargin=*]\item PySUS Documentation: \url{https://pysus.readthedocs.io/en/latest/}\item PySUS Data Sources: \url{https://pysus.readthedocs.io/en/latest/databases/data-sources.html}\item Pipeline NSS: \url{https://github.com/Zadoque/NSS-pipeline}\item Frontend NSS: \url{https://github.com/Zadoque/nss-front-end}\item Deployment NSS: \url{https://github.com/Zadoque/NSS-deployment}\item NixOS: \url{https://github.com/Zadoque/nss-NixOS-Configuration}\end{itemize}
\end{document}
